\chapter{Exam Mapping: AWS Certified Data Engineer -- Associate}
\index{certification!exam mapping}\index{AWS Certified Data Engineer}%
This appendix maps the AWS Certified Data Engineer -- Associate exam domains and objectives onto
the book's chapters and sections \cite{awsCertDataEngineer}, then gives the study strategy and a
self-assessment checklist. Weights are indicative of the published exam guide; always confirm
against the current guide before booking.

\section{Domain-to-Chapter Map}
\index{exam domains}%

\begin{table}[H]
\centering
\small
\begin{tabular}{@{}p{4.6cm}p{1.8cm}p{6.6cm}@{}}
\toprule
\textbf{Exam domain} & \textbf{Weight} & \textbf{Book coverage} \\
\midrule
1. Data Ingestion and Transformation & $\sim$34\% & Ch.\ 1 (S3 events, transfer), 9 (EMR), 10 (Glue ETL), 11 (DMS CDC, contracts), 13 (Kinesis), 14 (Firehose, MSK), 15 (Flink), 20 (Textract/Comprehend) \\
2. Data Store Management & $\sim$26\% & Ch.\ 1 (S3 classes, lifecycle), 2 (RDS/Aurora), 3 (DynamoDB), 4 (ElastiCache/DAX), 5--6 (Redshift), 7 (lakehouse, Athena, Spectrum) \\
3. Data Operations and Support & $\sim$22\% & Ch.\ 12 (MWAA, Step Functions), 16 (serving, monitoring), 19 (pipelines, drift), 22 (data quality), 24 (DataOps, CI/CD) \\
4. Data Security and Governance & $\sim$18\% & Ch.\ 8 (Lake Formation, sharing), 21 (IAM, VPC endpoints, KMS), 22 (lineage, mesh), 23 (Macie, erasure) \\
\bottomrule
\end{tabular}
\caption{Exam domains mapped to the book (same mapping as Chapter~24, expanded below).}
\label{tab:app-exam-domains}
\end{table}

\section{Objective-Level Breakdown}
\subsection{Domain 1: Data Ingestion and Transformation}
\begin{itemize}
    \item \textbf{Ingest data from sources:} Kinesis producer/shard design (Ch.~13, Monday--Wednesday);
    Firehose delivery, transformation, backup (Ch.~14); DMS full-load-plus-CDC (Ch.~11); S3 event
    notification patterns (Ch.~1 Thursday, Ch.~20).
    \item \textbf{Transform and process data:} Glue Studio and PySpark jobs, DynamicFrames,
    bookmarks (Ch.~10); EMR deployment modes (Ch.~9); Flink SQL windows, watermarks, exactly-once
    (Ch.~15); Lambda event processing (Ch.~14--15, 20).
    \item \textbf{Orchestrate pipelines:} MWAA DAGs, sensors, retries, trigger rules; Step
    Functions state machines (Ch.~12).
\end{itemize}
\subsection{Domain 2: Data Store Management}
\begin{itemize}
    \item \textbf{Choose a data store:} the decision framing of Ch.~2 (relational), Ch.~3 (key-value),
    Ch.~4 (cache), Ch.~5 (warehouse), Ch.~7 (lakehouse)---the exam's ``which service for these
    constraints'' questions live here.
    \item \textbf{Design data models and schemas:} DynamoDB single-table design and GSIs (Ch.~3);
    Kimball grain, fact types, SCD (Ch.~5); distribution and sort keys (Ch.~6); schema evolution
    and compatibility (Ch.~11).
    \item \textbf{Manage the data lifecycle:} S3 storage classes and lifecycle rules (Ch.~1);
    retention, versioning, time travel (Ch.~1, 7); TTL semantics (Ch.~3).
\end{itemize}
\subsection{Domain 3: Data Operations and Support}
\begin{itemize}
    \item \textbf{Automate processing:} idempotent re-runs, bookmarks, manifests (Ch.~10--12);
    IaC with Terraform/CDK and CI/CD gates (Ch.~24).
    \item \textbf{Monitor and troubleshoot:} CloudWatch metrics that matter---CDC lag
    (Ch.~11), \texttt{IteratorAgeMilliseconds} (Ch.~13), WLM queue wait (Ch.~6), delivery success
    (Ch.~14); pipeline failure alerting (Ch.~12, 22).
    \item \textbf{Maintain data quality:} Glue Data Quality rulesets and thresholds (Ch.~22);
    contract enforcement (Ch.~11); dbt tests (Ch.~24).
\end{itemize}
\subsection{Domain 4: Data Security and Governance}
\begin{itemize}
    \item \textbf{Authentication and authorization:} IAM roles, policies, boundaries, SCPs
    (Ch.~21); Lake Formation grants, LF-Tags, column-level security (Ch.~8); QuickSight RLS
    (Ch.~16).
    \item \textbf{Encryption and data protection:} KMS key policies, Secrets Manager rotation
    (Ch.~21); S3 Object Lock and SSE options (Ch.~1); crypto-shredding (Ch.~23).
    \item \textbf{Logging, auditing, and governance:} CloudTrail including data events (Ch.~23);
    catalog and lineage (Ch.~7, 22); Macie discovery and quarantine (Ch.~23); clean rooms and
    data sharing (Ch.~8).
\end{itemize}

\section{Study Strategy}
\index{study strategy}%
\begin{enumerate}
    \item \textbf{Self-assess per domain.} For each domain, re-read the chapter Key Takeaways and
    attempt the exercises cold. Mark each objective \emph{confident}, \emph{shaky}, or
    \emph{blank}.
    \item \textbf{Repair blanks by doing.} A shaky area you have used recovers by re-reading; a
    blank area recovers only by running the lab (Appendix~D has the reference solutions).
    \item \textbf{Practice scenario questions.} The exam's character is ``which service/design
    for this constraint set''---the Friday use cases train exactly that muscle.
    \item \textbf{Two full-length mock exams,} spaced two days apart, in the final week; review
    every wrong answer back to its chapter. Target: 80\%+ overall, no domain below 70\%.
    \item \textbf{Use the cross-cloud appendix defensively.} Questions sometimes name a
    non-AWS concept (Kafka, Airflow, Delta); knowing the AWS equivalent (MSK, MWAA, Glue/EMR)
    converts them into free points.
\end{enumerate}

\begin{tipbox}
Exam heuristics that recur across domains: prefer managed/serverless when the stem emphasizes
operational simplicity; prefer provisioned when it emphasizes steady-state cost or fine control;
prefer the answer that encrypts, logs, and least-privileges by default; distrust any answer whose
durability story is ``a job remembers to clean up.''
\end{tipbox}

\section{Self-Assessment Checklist}
\index{self-assessment}%
\providecommand{\chk}{\item[$\square$]}
Tick an item only when you can do it unaided---build it or explain it---not merely recognize it.

\subsection*{Domain 1: Ingestion and Transformation}
\begin{itemize}
\chk Compute Kinesis shard counts from byte and record rates; choose partition keys; interpret
\texttt{IteratorAgeMilliseconds}.
\chk Build a Firehose stream with Lambda transform, Parquet conversion, and raw backup.
\chk Configure DMS full-load-plus-CDC and read its lag metrics.
\chk Write a Glue job with bookmarks and explain what bookmarks do \emph{not} guarantee.
\chk Write Flink SQL with event-time windows and a declared watermark.
\chk Author an MWAA DAG with sensors, retries, and a failure alert that cannot be skipped.
\end{itemize}
\subsection*{Domain 2: Data Store Management}
\begin{itemize}
\chk Choose between S3 classes and write a lifecycle policy for a stated retention requirement.
\chk Design a DynamoDB table from an access-pattern list, including GSIs and TTL.
\chk Explain Multi-AZ versus read replicas, and when Aurora Serverless v2 wins.
\chk Declare a fact-table grain; choose distribution and sort keys; read an \texttt{EXPLAIN} plan.
\chk Explain why Parquet + partitioning collapses Athena bytes scanned.
\end{itemize}
\subsection*{Domain 3: Operations and Support}
\begin{itemize}
\chk Make any pipeline in this book re-runnable without duplicates and say exactly how.
\chk Wire a data-quality gate that halts a pipeline and pages a human.
\chk Read a failed Step Functions execution and locate the failing state.
\chk Explain step caching in SageMaker Pipelines and when it masks a real change.
\end{itemize}
\subsection*{Domain 4: Security and Governance}
\begin{itemize}
\chk Revoke \texttt{IAMAllowedPrincipals} and grant column-level access with Lake Formation.
\chk Design private S3 access with no internet route and an endpoint policy.
\chk Explain why KMS key policies complete the network perimeter.
\chk Trace a bad batch to its consumers using lineage.
\chk Design an erasure mechanism (crypto-shredding or enumerated deletion) for a stated
regulation.
\end{itemize}
